% ID: SEL-D0053-OA-E02
% ESTADO: cotejada. Segunda pasada visual y borrador compilado revisados.
% Fuente: Selectividad/sel_2013_matematicas_aplicadas/septiembre/titular/EXAMEN 2-ANDALUCÍA.pdf
% SHA-256: f8633bbd91d33bee02c2cc582c33b7ffeb3643f312b810ad0bb8c5d154d0721d
% Página: 1; opción A; ejercicio original 2.

En una empresa de montajes el número de montajes diarios realizados por un trabajador depende de los días trabajados según la función $M(t)=\dfrac{11t+17}{2t+12}$, $t\geq1$, donde $t$ es el número de días trabajados.
\begin{enumerate}[label=\alph*)]
\item \textbf{(0.5 puntos)} \textquestiondown{}Cuántos montajes realiza el primer día? \textquestiondown{}Cuántos días necesitará para realizar cinco montajes diarios?
\item \textbf{(0.75 puntos)} \textquestiondown{}Qué ocurriría con el número de montajes diarios si trabajara indefinidamente?
\item \textbf{(0.75 puntos)} El dueño de la empresa cree que el número de montajes diarios aumenta con los días de trabajo. Estudiando la función, justifique si es cierta dicha creencia.
\item \textbf{(0.5 puntos)} Dibuje la gráfica de la función.
\end{enumerate}
